
\newcommand{\CaptionAnatomy}{\textbf{Benchmark anatomy.} Evaluations per (A) functional system, (B) process class, (C) Biological level, (D) workflow stage, and (E) data type. Light bars in A and B are sub-categories: Cross-system is split by tissue compartment or readout, and process class is split into normal function and disease, injury and aging, then into sub-categories and disease classes. Each evaluation has one label on each taxonomy axis; an evaluation can use more than one data type, so E sums to more than \NumEvals{}.}
\newcommand{\CaptionLeaderboard}{\textbf{Leaderboard and run characteristics of 34 model--harness configurations.} (A) Pass rate of each configuration over 91 evaluations and three attempts each, ordered by rank; whiskers are 95\% Student's t intervals over per-evaluation pass rates and colors mark the model provider. (B--E) Pass rate against the mean cost (B), tokens (C), wall-clock time (D) and number of tool calls (E) per run, on log axes; marker shape gives the harness, numbers give the rank in A for the Pareto front and the top three configurations, and the dashed line traces the front. Cost was not reported for five of the six Codex configurations and tool calls for none of them, so those points are omitted from B and E. (F) For each configuration, the share of evaluations passed in every attempt, in some attempts, or never. (G) Pass rate of the 13 models run under two harnesses; the line joins the two configurations of a model. Refusals (4 runs) and platform failures (66 runs) are excluded throughout, leaving 9,212 graded attempts.}
\newcommand{\CaptionStages}{\textbf{Pass rate by data type.} Mean pass rate per data type, grouped by the biological scale of the data; the dashed line is the mean over the \NumModels{} models. An evaluation can use more than one data type; n is the number of evaluations using each. One flagship model from each lab in color and every other model in grey; both harnesses pooled and averaged at the evaluation level.}
\newcommand{\CaptionMatrix}{\textbf{Pass rate by biological level, process class, and functional system.} All \NumModels{} models, both harnesses pooled, ordered by overall pass rate. Each evaluation carries at most one biological level, process class, and functional system classification. The outlined cell is the leader of each column and the bottom row is the mean over models.}
\newcommand{\CaptionTraps}{\textbf{Mean pass rate by analytical trap and inference fallacy.} (A) Analytical traps. (B) Inference fallacies. Bars give the mean over the 21 models of each model's pass rate on the evaluations in the row, both harnesses pooled, ordered from hardest to easiest; n is the number of evaluations.}
